%% Geometric Encryption: Property-Preserving Database Encryption
%% via Gauge Invariance on Fiber Bundles
%%
%% Bee Rosa Davis · Davis Geometric · 2026
%% Companion paper to the Kähler Substrate (Davis 2026).
%% Patent reference: U.S. Provisional Application No. 64/045,889 (GIGI),
%% with continuation-in-part covering v0.3 constructions to be filed
%% before submission.

\documentclass[11pt,letterpaper]{article}
\usepackage[utf8]{inputenc}
\usepackage[T1]{fontenc}
\usepackage{lmodern}
\usepackage{textcomp}
\usepackage[margin=1in]{geometry}
\emergencystretch=3em
\usepackage{amsmath,amssymb,amsthm}
\usepackage{mathtools}
\usepackage{hyperref}
\usepackage{booktabs}
\usepackage{multirow}
\usepackage{float}
\usepackage{tabularx}
\usepackage{needspace}
\usepackage{graphicx}
\usepackage{enumitem}
\usepackage{listings}
\usepackage{xcolor}

\newtheorem{theorem}{Theorem}[section]
\newtheorem{lemma}[theorem]{Lemma}
\newtheorem{proposition}[theorem]{Proposition}
\newtheorem{corollary}[theorem]{Corollary}
\newtheorem{definition}[theorem]{Definition}
\newtheorem{remark}[theorem]{Remark}

\hypersetup{
  colorlinks=true,
  linkcolor=black,
  urlcolor=blue,
  citecolor=black,
}

\title{\textbf{Geometric Encryption} \\[6pt]
\large Property-Preserving Database Encryption \\
via Gauge Invariance on Fiber Bundles}

\author{Bee Rosa Davis \\
\href{mailto:bee_davis@alumni.brown.edu}{\nolinkurl{bee_davis@alumni.brown.edu}} \\
Davis Geometric, Sacramento, CA}

\date{\today}

\begin{document}
\maketitle

\begin{abstract}
\noindent We introduce \emph{geometric encryption}, a property-preserving
database-encryption framework in which the encryption acts on the
fiber of a fiber-bundle data store. The action determines, by
construction, which queries are computable on ciphertext: a query
of arity $m$ is computable on ciphertext at native server speed
with $O(\mathrm{poly}(m))$ closed-form client post-processing iff
it is \emph{$\rho$-equivariant} under the structure group $G$ for
some representation $\rho: G \to \mathrm{Aut}(\mathbb{R}^m)$. The
$\rho = \mathrm{id}$ slice (gauge-invariant queries) needs no
post-processing; the non-trivial slices --- SUM, AVG, MIN, MAX,
VAR, STDDEV, RANGE on affine ciphertext --- are recovered by the
client via a single closed-form $\rho_g^{-1}$ application. We propose a taxonomy of five encryption modes (Affine,
Opaque, Indexed, Probabilistic, Isometric), each parameterized by its
mathematical structure --- a fiber-automorphism group for the
deterministic modes (Affine: $\mathrm{Aff}(\mathbb{R})^k$;
Isometric: $O(k)$ on a grouped $k$-tuple), and a randomized or
one-way structure map for the modes whose action is not a
deterministic bijection (Opaque: AEAD-randomized; Indexed: PRF
quotient; Probabilistic: $\mathrm{Aff}(\mathbb{R})$ with Gaussian-perturbed
key). On top of the taxonomy we develop five higher-level
constructions: Curvature-MAC bundle integrity,
$\mathrm{Aff}(\mathbb{R})$ capability delegation, holonomy ledger,
\v{C}ech threshold sharing, and a continuous RG-flow ratchet for
per-write forward secrecy. A sixth construction --- a
pairing-based collusion-resistant alternative to
$\mathrm{Aff}(\mathbb{R})$ delegation, built on BLS12-381 with
collusion resistance reducing to discrete log on $G_2$
(${\sim}2^{128}$ work) --- ships alongside. All five shipping modes
and six constructions are implemented in the public GIGI Rust
codebase~\cite{gigi2026code} with $981+$ passing tests; an
independent Python math-validation suite ($66$ tests, including
the $66/66$ FHE / PQ parity rigor oracle that proves the
order-statistic Probabilistic-bias caveat) confirms every
shipped mathematical claim from a second implementation path. We
document one resolved floating-point design issue: the v0.3.0 draft
included Davis capacity $\tau / K$ in the integrity-tag input, which
amplified $K$'s relative drift through the division and required a
$10^{-6}$ quantization grain to absorb. The v0.3.1 design replaces
the capacity slot with the record-count $\tau$ (an exact $u64$) and
tightens the f64 quantization grain to $10^{-10}$. Three further
constructions (geodesic-ball ABE, spectral-signature ZKP,
lattice-fiber post-quantum gauge) are stated as derived but with
reference implementation deferred to a successor paper. The
$\mathrm{Aff}(\mathbb{R})$ delegation construction is explicitly
\emph{not} collusion-resistant proxy re-encryption in the
Ateniese--Hohenberger sense; this is a fundamental consequence of
additive arithmetic, documented as a tested limitation alongside the
collusion-resistant pairing alternative. \emph{Draft status}: this
paper is v0.1 and presents \S\S 1--3 (introduction, background,
gauge-invariance framework). \S\S 4--14 (full mode taxonomy,
constructions, evaluation) are in active draft; the test counts and
performance numbers above are reproducible from the public
repository today.
\end{abstract}

\tableofcontents
\bigskip


%% ==================================================================
\section{Introduction}
\label{sec:intro}
%% ==================================================================

\subsection{The compute-on-ciphertext problem}
\label{sec:intro:problem}

Three decades after Rivest, Adleman, and Dertouzos posed it
in 1978, the question ``how do we compute on encrypted data?'' remains
the central tension of database confidentiality. The dominant answers
cluster into two families with sharp trade-offs:

\begin{enumerate}[nosep]
\item \textbf{Fully homomorphic encryption} (Gentry 2009 and
  successors~\cite{gentry2009fhe,bgv2012,ckks2017}) provides universality
  at a $10^3$--$10^5\times$ slowdown. Any function on plaintext can be
  evaluated on ciphertext, at a cost that has slowly come down but
  remains incompatible with production latency budgets for analytical
  queries.
\item \textbf{Transparent data encryption} (Oracle TDE, AWS RDS
  encryption, MongoDB CSFLE) provides at-rest confidentiality at a
  $1\times$ cost --- but requires decryption to compute. Every
  analytical query passes through plaintext at some point on a
  privileged node.
\end{enumerate}

Between these poles sits a smaller family of \emph{property-preserving
encryption} (PPE) constructions: order-preserving
encryption~\cite{boldyreva2009ope,popa2011cryptdb}, structured
encryption~\cite{chase2010structured}, searchable symmetric
encryption~\cite{curtmola2006sse}, and functional
encryption~\cite{boneh2011fe}. Each preserves one specific property
(order, equality, a chosen function class). None offer a unifying
framework that says ``encryption mode $X$ preserves query class $Y$
because of mathematical reason $Z$.'' We provide such a framework.
The closest spiritual antecedent in the cryptographic literature is
\emph{privacy-preserving blueprints}~\cite{kohlweiss2023blueprints},
which prove that committed values satisfy a stated policy without
revealing the values themselves; the gauge-invariant ring developed
in \S\ref{sec:gauge-invariance} is the database-runtime
instantiation of the same idea, with the structure group of the
encryption playing the role of the policy.

\subsection{The geometric move}
\label{sec:intro:geometric}

The Davis framework~\cite{davis2026sheafcomposition,davis2026manifold}
represents a database as sections of a fiber bundle
$(E, B, F, \pi, \Phi)$, where $B$ is the base space (keys), $F$ is the
fiber (non-key field values), and $\pi: E \to B$ is the projection. A
record is a section $\sigma: p \mapsto (p, v_1, \ldots, v_k) \in E$.

The framework's analytics --- curvature, anomaly detection, spectral
connectivity, prediction --- are functions on the bundle's intrinsic
geometry. Crucially, they are \emph{invariants} of the bundle's
geometric structure: they do not depend on which coordinate system the
fiber is expressed in.

Geometric encryption is the observation that the choice of coordinate
system on the fiber \emph{is} the encryption. A secret gauge
transformation $g \in G \subset \mathrm{Aut}(F)$ scrambles each
record's fiber values $(v_1, \ldots, v_k) \mapsto g(v_1, \ldots, v_k)$.
An attacker who sees only the gauged data cannot recover plaintext
(without knowing $g$); but anyone --- including analysts without
decryption privilege --- can compute any $G$-invariant function on the
gauged data and get the same answer as on plaintext.

The encryption is not a layer bolted onto the database. \textbf{The
geometry is the encryption.}

\subsection{What this paper adds}
\label{sec:intro:adds}

\begin{enumerate}[nosep]
\item \textbf{A closed taxonomy} of five encryption modes (Affine,
  Opaque, Indexed, Probabilistic, Isometric) keyed by their structure
  groups. Each row of the taxonomy is a triple (structure group $G$,
  AEAD/PRF primitive, query class preserved).
\item \textbf{Six higher-level constructions} built on the taxonomy
  (five in the main path plus a pairing-based collusion-resistant
  delegation alternative):
  Curvature-MAC bundle integrity, $\mathrm{Aff}(\mathbb{R})$ capability
  delegation, holonomy ledger (Merkle audit log), \v{C}ech threshold
  sharing, continuous RG-flow ratchet.
\item \textbf{A reference implementation} in the public GIGI Rust
  codebase~\cite{gigi2026code}, with TDD test surface ($826+$ passing
  tests) backing every claim. Each math primitive ships with both Rust
  integration tests and Python math-validation mirrors.
\item \textbf{Honest scope statements} for each construction's
  threat model. $\mathrm{Aff}(\mathbb{R})$ capability delegation
  (\S 6) is \emph{not} collusion-resistant PRE in the
  Ateniese--Hohenberger sense~\cite{ateniese2005pre,libert2008pre} ---
  a fundamental consequence of additive arithmetic. The
  $(\alpha, \beta)$ capability plus the delegatee's own gauge key
  yields a 2-equations-in-2-unknowns linear system whose solution
  is the delegator's secret in $O(1)$ field operations. This means
  Aff$(\mathbb{R})$ delegation requires \emph{both} the delegatee
  being honest \emph{and} the delegatee's key remaining
  uncompromised; a key-leak in the delegatee compounds with a
  capability leak to immediately expose the delegator.
\item For deployments where this combined-trust model is too
  strong, we ship a pairing-based \emph{collusion-resistant
  key-delegation} primitive in
  \texttt{src/pairing\_delegation.rs}~\cite{gigi2026code} (related
  in spirit to the production PRE of
  Umbral~\cite{nunez2018umbral} but simpler in shape). The
  construction is conceptually adjacent to two decades of
  anonymous-credential delegation work originating with the CL
  signature scheme~\cite{camenisch2004clsigs} and continuing through
  mercurial signatures~\cite{camenisch2019mercurial,
  garman2024mercurial}; we deliver the database-record-access
  analogue of what those constructions deliver at the credential
  layer. Construction:
  BLS12-381 keypairs $sk \in \mathbb{F}_p$, $pk = g_2^{sk} \in G_2$;
  KEM capsule $C_1 = g_1^r \in G_1$; delegation capability
  $rk_{A \to B} = pk_B^{1/sk_A} = g_2^{sk_B / sk_A} \in G_2$. We
  emphasize \emph{this is not strict proxy re-encryption}: the
  proxy holds $rk$ and forwards it but plays no cryptographic
  role; the delegatee applies the capability themselves via
  pairing,
  \[
  K_{\mathrm{shared}} \;=\; e(C_1,\, rk_{A \to B})
                       \;=\; e(g_1, g_2)^{r \cdot sk_B / sk_A}.
  \]
  Alice arrives at the same $K_{\mathrm{shared}}$ via her own
  side-computation (she knows $sk_A$ and $pk_B$, so she can
  construct $rk$ and apply the pairing herself). The two parties'
  AEAD keys are derived from the same $K_{\mathrm{shared}}$ via
  HKDF-SHA256, enabling sender-receiver agreement. The full
  Ateniese--Hohenberger AHP construction with a separate
  $g_1^r$-bound commitment encrypted under the delegatee's pubkey
  --- which would make the proxy's operation cryptographic in the
  strict sense --- is left to a v0.4 refinement.
  \par
  \par
  \textbf{Security (informal statement; full proof in \S 6).}
  Let $\mathcal{A}$ be a PPT adversary in the trusted-delegatee
  model given the public parameters $(G_1, G_2, G_T, e, g_1, g_2)$
  for BLS12-381, Alice's public key
  $pk_A = g_2^{sk_A}$, the delegation capability
  $rk_{A \to B} = g_2^{sk_B / sk_A}$, the KEM capsule $C_1 = g_1^r$
  for an honestly-sampled $r \in \mathbb{F}_p$, and oracle access
  to all algorithms on non-challenge inputs. Then $\mathcal{A}$'s
  distinguishing advantage against the delegated-KEM shared key
  $K_{\mathrm{shared}} = e(g_1, g_2)^{r \cdot sk_B / sk_A} \in G_T$
  reduces to the \emph{Bilinear Diffie--Hellman} (BDH) assumption
  on BLS12-381:
  \[
  \mathrm{Adv}^{\mathrm{deleg\text{-}KEM}}_{\mathcal{A}}
  \;\leq\;
  \mathrm{Adv}^{\mathrm{BDH}}_{\mathcal{B}},
  \]
  for a constructed PPT $\mathcal{B}$ running in the same time as
  $\mathcal{A}$. Reduction sketch: embed the BDH challenge
  $(g_1, g_1^a, g_1^b, g_2, g_2^a, g_2^b)$ by setting
  $C_1 = g_1^a$ and $rk_{A \to B} = g_2^b$ (implicit
  $r = a$, $sk_B / sk_A = b$); then $K_{\mathrm{shared}}
  = e(g_1, g_2)^{ab}$ is exactly the BDH target. The
  \emph{collusion} case (delegatee $\cup$ proxy) reduces to DLP
  on $G_2$: the coalition computes
  $rk^{1/sk_B} = g_2^{1/sk_A} \in G_2$, and recovering
  $sk_A \in \mathbb{F}_p$ from $g_2^{1/sk_A}$ requires solving the
  discrete logarithm on $G_2$, hard at $\sim 2^{128}$ work in
  BLS12-381. Forging $K_{\mathrm{shared}}$ for a fresh capsule
  $C_1' = g_1^{r'}$ unseen by the coalition reduces to DLP on
  $G_1$ (recovering $r'$) or to BDH (computing the pairing
  output directly). The construction therefore upgrades the
  Aff$(\mathbb{R})$ algebraic-recovery attack (constant-time) to
  a computational-hardness floor under BDH (passive adversary)
  and DLP (active collusion). Both constructions share the paper,
  with \S 6 documenting the Aff$(\mathbb{R})$ trusted-delegatee
  model and the pairing construction's BDH-hardness theorem
  side-by-side; the formal theorem statement and complete
  reduction live in \S 6.
\item \textbf{Three derived-but-unimplemented constructions}
  (geodesic-ball ABE, spectral-signature ZKP, lattice-fiber PQ gauge)
  with security arguments, deferred to a v0.4 successor paper. The
  Davis-Manifold pattern: publish the construction, defer the production
  implementation until peer review of the construction.
\end{enumerate}

\subsection{Scope: what this paper is and is not}
\label{sec:intro:not}

\paragraph{Relation to FHE.}
GIGI is a different point on the Pareto frontier from fully
homomorphic encryption~\cite{gentry2009fhe,bgv2012,ckks2017}. FHE
achieves arbitrary-function evaluation on ciphertext at
$10^3$--$10^5\times$ overhead; GIGI achieves a large
\emph{gauge-equivariant} query class at native speed via $O(1)$
client-side post-processing per query. Concretely, the
gauge-equivariant class shipped in v0.3 (\texttt{src/aggregate\_helpers.rs})
covers:
\begin{quote}
\sloppy
COUNT, SUM, AVG, MIN, MAX, VARIANCE, STDDEV, RANGE, any polynomial
combination of these, GROUP BY with these aggregates, HAVING on
aggregate predicates, range queries via literal-transform, equality
search (via INDEXED mode), distance / k-NN (via ISOMETRIC mode),
$\sigma$-bucket approximate equality (via PROBABILISTIC mode).
\end{quote}
The pattern: server computes $f(\text{ciphertext}) = h(f(\text{plaintext}), g)$
for closed-form $h$; client applies $h^{-1}$ to a single scalar.
The record set never leaves ciphertext form on the server.

\paragraph{Caveat: order statistics under Probabilistic mode.}
The claim above is \emph{plaintext-exact} for COUNT, SUM, AVG, VAR,
STDDEV under both Affine and Probabilistic modes (verified in
\texttt{tests/fhe\_pq\_parity\_rigor.rs} \S\,A.1--A.3 and
\texttt{validation\_tests\_fhe\_pq\_rigor.py} \S\,A.1--A.3), and for
MIN, MAX, RANGE under Affine mode. \emph{Under Probabilistic mode with
$\sigma > 0$, MIN, MAX, RANGE are biased}: order statistics do not
commute with additive Gaussian noise. The naïve closed-form inverse
$(f_{\text{cipher}} - b) / a$ overshoots by
$\Theta(\sigma \sqrt{2 \log n} / |a|)$ in expectation
(sub-Gaussian extreme-value bound), and the bias does \emph{not}
vanish as $n \to \infty$. The shipped helpers
\texttt{decrypt\_min}, \texttt{decrypt\_max}, \texttt{decrypt\_range}
return \texttt{BiasedUnderProbabilisticNoise} when $\sigma > 0$;
callers who accept the bias for coarse bounds use the explicit
\texttt{*\_unchecked} variants. Bias detectability is locked in by
\texttt{tests/fhe\_pq\_parity\_rigor.rs::a4\_*\_is\_biased\_*} and the
Python oracle \S\,A.4. Practical mitigations: (i) use VAR + mean for
spread queries (VAR \emph{is} bias-correctable via $\sigma^2$
subtraction); (ii) switch the column to Affine mode (set $\sigma = 0$)
for exact MIN/MAX; (iii) pair GIGI with FHE on the specific column
when exact encrypted-side extremes under noise are required.

The genuine FHE-only territory is narrow: encrypted-side
conditionals whose branch depends on encrypted predicates,
non-polynomial functions of fiber values (transcendentals,
neural-network activations), and reductions whose intermediate
values must stay encrypted (median, $\arg\max$). For these we
recommend pairing GIGI with an FHE layer on the specific columns
that need them, rather than wrapping the entire database in FHE.

\paragraph{Confidentiality scope.}
The framework provides \emph{selective compute on encrypted data}.
For confidentiality of individual record values --- ensuring an
attacker who reads the database file cannot recover plaintext ---
the OPAQUE mode (AES-256-GCM-SIV) is the appropriate primitive,
applied to columns containing sensitive values. The
gauge-equivariant modes (Affine, Probabilistic, Isometric) are
not chosen-plaintext-attack secure on their own; in deployments
where an attacker may obtain plaintext--ciphertext pairs, use
OPAQUE for those columns.

\paragraph{Quantum security.}
The taxonomy of shipped primitives by post-quantum status:

\begin{center}
\small
\begin{tabular}{p{4.5cm}p{6.5cm}c}
\toprule
\textbf{Primitive} & \textbf{Used for} & \textbf{PQ?} \\
\midrule
$\mathrm{Aff}(\mathbb{R})$, $\mathrm{Probabilistic}$, $O(k)$ & gauge modes (linear arithmetic, no algebraic hardness assumption) & \checkmark \\
AES-256-GCM-SIV & OPAQUE mode AEAD & \checkmark \\
AES-256-CMAC & INDEXED mode PRF & \checkmark \\
HMAC-SHA256 & Curvature-MAC, Čech share authentication & \checkmark \\
HKDF-SHA256 & integrity key derivation, ratchet chain & \checkmark \\
SHA-256 & holonomy ledger Merkle leaves and internal nodes & \checkmark \\
Shamir over secp256k1 base field $\mathbb{F}_p$ & Čech threshold sharing (we use the \emph{field}, not the curve group) & \checkmark \\
ML-KEM-768 (FIPS 203) & PQ trusted-delegatee mode (Sprint J.3), threshold-PQ delegation envelopes (Sprint J.4) & \checkmark \\
BLS12-381 pairings ($G_1, G_2, G_T$) & collusion-resistant delegation (Sprint J.2; pre-quantum, retained alongside the two PQ-safe delegation modes below) & $\times$ \\
\bottomrule
\end{tabular}
\end{center}

Of the ten shipped primitive classes, nine are post-quantum
acceptable at NIST Level 1 or higher (AES-256 family, SHA-256-based
constructions, ML-KEM-768 at Level 3). The exception is the
BLS12-381 pairing-based delegation, whose security reduces to BDH
on $G_T$ and DLP on $G_2$ --- both broken by Shor's algorithm.
\textbf{For deployments requiring post-quantum security
across all delegation paths, v0.3 ships two PQ-safe alternatives}
that together cover both threat models:
\begin{itemize}[nosep]
\item ML-KEM trusted-delegatee delegation (Sprint J.3, NIST Level 3
  PQ under MLWE) --- single-party recipient, no collusion-resistance
  claim; the analogue of the BLS12-381 mode under the PQ assumption.
\item Lattice threshold delegation (Sprint J.4, Shamir K-of-N over
  $\mathbb{F}_p$ composed with per-share ML-KEM transport).
  Information-theoretic collusion-resistance against any
  $K-1$-shareholder coalition (Shamir) composed with PQ transport
  (MLWE) --- a \emph{strictly stronger} collusion-resistance
  property than BLS12-381's computational DLP-based resistance, but
  predicated on the multi-party deployment shape.
\end{itemize}
The pre-quantum BLS12-381 mode is retained for deployments that
need (i) a single-party recipient and (ii) collusion-resistance
under classical assumptions; PQ-safe single-party
collusion-resistant delegation (lattice-PRE in the Kirshanova~2014
sense) is the v0.4 deferred construction. Users select the mode
matching their threat model.

The full post-quantum table: GIGI v0.3 ships PQ-safe primitives
for confidentiality (OPAQUE), equality search (INDEXED),
integrity (Curvature-MAC + Holonomy Ledger), key derivation
(HKDF), threshold sharing (Čech), continuous forward secrecy
(Sprint M ratchet via HKDF), and two flavors of delegation
(ML-KEM trusted + lattice threshold).

\paragraph{Threshold vs.\ pairing-PRE: trust-model dependence.}
The lattice threshold construction (Sprint J.4) and the BLS12-381
pairing PRE (Sprint J.2) are \emph{not} substitutable across all
trust models; the right choice depends on the deployment:

\begin{itemize}[nosep,leftmargin=*]
\item For \textbf{$K$-of-$N$ distributed delegation} (Alice
  delegates to a quorum), the lattice threshold construction is
  strictly stronger than pairing PRE: it is post-quantum (ML-KEM-768
  outer + Shamir over $\mathbb{F}_p$ inner) and any subset of size
  $\leq K-1$ recovers \emph{information-theoretically zero} ---
  exhaustively verified for $(K{=}3, N{=}5)$ across all $\binom{5}{2}
  = 10$ collusion subsets in
  \texttt{tests/fhe\_pq\_parity\_rigor.rs::b2\_*\_subset\_*}. This
  is stronger than the DLP-on-$G_2$ hardness pairing PRE relies on,
  and dominates pairing PRE on the multi-party axis.
\item For \textbf{single-party delegation Alice $\to$ single Bob}
  with delegatee-collusion resistance, the threshold construction
  degenerates at $K = 1$ to plain ML-KEM trusted delegation: no
  collusion-resistance beyond ``Bob is trusted.'' Pairing PRE, by
  contrast, gives delegatee-vs-proxy collusion-resistance via
  $\mathrm{DLP}_{G_2}$-hardness (Sprint J.2) at the cost of being
  pre-quantum. True single-party post-quantum collusion-resistant
  PRE (lattice-PRE in the Kirshanova~2014 / Aono--Hayashi~2017
  sense) is the v0.4 deferred construction.
\end{itemize}

Therefore: \textbf{lattice threshold dominates pairing PRE on the
quorum axis but does not replace it on the single-delegatee axis};
both are shipped and users select per their trust model.

\subsection{Roadmap}
\label{sec:intro:roadmap}

\S\ref{sec:background} reviews the prior canon --- three Davis
manuscripts that this work extends --- plus the external math anchors.
\S\ref{sec:gauge-invariance} states the gauge-equivariance
theorem: a query is computable on encrypted data with at most
$O(\mathrm{poly}(m))$ client-side post-processing iff it is
$\rho$-equivariant under the structure group $G$ for some
representation $\rho: G \to \mathrm{Aut}(\mathbb{R}^m)$. The
invariant case ($\rho = \mathrm{id}$) is the no-post-processing
slice. \S\S 4--9 (forthcoming) develop the five
modes and six constructions in turn. \S 10 reports the empirical
evaluation; \S 11 presents the three deferred constructions; \S 12
surveys related work; \S 13 limits and future work; \S 14 conclusion.


%% ==================================================================
\section{Background}
\label{sec:background}
%% ==================================================================

\subsection{Prior canon --- three Davis manuscripts}
\label{sec:background:canon}

This paper is the cryptographic application of two prior Davis
substrates:

\paragraph{The Davis Manifold~\cite{davis2026manifold}.}
A Riemannian substrate $(M, g, P, \varepsilon, \kappa_{\text{hard}},
\kappa_{\text{soft}})$ for residue-bounded computation. The non-vacuity
condition $\kappa_{\text{soft}} - 2R\varepsilon(L^*) > 0$ is the source
of the Hadamard-detection invariant that Sprint M (continuous ratchet,
\S 9) preserves under per-write gauge advance.

\paragraph{The K\"ahler Substrate~\cite{davis2026kahlersubstrate}.}
Extends $(M, g)$ to $(M, g, J, \nabla, B, \Gamma)$ --- a K\"ahler
manifold with closed magnetic 2-form $B$ and discrete graph
approximation $\Gamma$. The eight-layer catalog (L1--L8 plus
extensions) provides the invariant ring (curvature, spectral gap,
Davis capacity, mean holonomy, Betti numbers) that Sprint I
(Curvature-MAC, \S 5) signs.

\paragraph{Sheaf Composition~\cite{davis2026sheafcomposition}.}
The discrete section-graph runtime --- the fiber-bundle data model on
which both substrates execute.

\subsection{External math anchors}
\label{sec:background:anchors}

\begin{itemize}[nosep]
\item \textbf{Property-preserving encryption}~\cite{pandey2012ppe}: the
  closest prior framework. PPE generalizes deterministic / OPE
  encryption to ``encryption that preserves a chosen property $P$.''
  Geometric encryption is PPE where $P$ is ``a chosen $G$-invariant
  geometric quantity'' and $G$ is the structure group of the encryption
  mode.
\item \textbf{Structured encryption}~\cite{chase2010structured}: framing
  for ``which queries are supported on which encrypted data
  structures.'' Closest in spirit but does not classify by structure
  group.
\item \textbf{Searchable symmetric encryption}~\cite{curtmola2006sse}:
  the search-on-ciphertext primitive that motivates the INDEXED mode of
  our taxonomy.
\item \textbf{Order-preserving encryption}~\cite{boldyreva2009ope,popa2011cryptdb}:
  the order-on-ciphertext primitive that the AFFINE mode of our
  taxonomy generalizes (Affine preserves order when $a > 0$; OPE is a
  special case).
\item \textbf{Homomorphic encryption}~\cite{gentry2009fhe,bgv2012,ckks2017}:
  the universality target. Different cost-vs-flexibility point.
\item \textbf{Proxy re-encryption}~\cite{ateniese2005pre,libert2008pre}:
  the standard collusion-resistant PRE definition. \S\ref{sec:intro:adds}
  noted explicitly that Sprint J does \emph{not} meet this; we
  document the gap in \S 6.
\item \textbf{AES-GCM-SIV}~\cite{rfc8452}: AEAD primitive for OPAQUE
  mode.
\item \textbf{AES-CMAC}~\cite{nist-sp800-38b}: PRF for INDEXED mode.
\item \textbf{HKDF}~\cite{krawczyk2010hkdf,rfc5869}: the KDF for
  INTEGRITY key derivation and Sprint M ratchet.
\item \textbf{Shamir secret sharing}~\cite{shamir1979}: the math
  underneath Sprint L.
\item \textbf{RFC 6962 Certificate Transparency
  log}~\cite{rfc6962}: the Merkle structure underneath Sprint K.
\item \textbf{Signal symmetric
  ratchet}~\cite{marlinspike2016signal}: the design analog for Sprint M
  continuous forward secrecy.
\item \textbf{CL signatures and mercurial
  signatures}~\cite{camenisch2004clsigs, camenisch2019mercurial,
  garman2024mercurial}: the two-decade anonymous-credential
  delegation lineage; conceptually adjacent to our pairing-based
  delegation construction (\S 6, Sprint J.2), which delivers the
  database-record-access analogue at the storage layer.
\item \textbf{Privacy-preserving
  blueprints}~\cite{kohlweiss2023blueprints}: the
  policy-on-committed-value proof primitive of which our
  gauge-invariant ring (\S\ref{sec:gauge-invariance:ring}) is the
  database-runtime instantiation.
\item \textbf{Verifiable random functions and standard-assumption
  PRFs}~\cite{lysyanskaya2002vrf, bellare2024prfs}: the PRF/VRF
  security foundations under our INDEXED mode (AES-256-CMAC).
\item \textbf{Dynamic accumulators}~\cite{camenisch2002accumulators,
  kemmoe2024accumulator}: the standard-cryptographic membership-proof
  primitive lineage; our field-index bitmap and base-graph topology
  provide a related geometric membership / spectral-connectivity
  primitive at the bundle layer.
\item \textbf{Authenticated data
  structures}~\cite{goodrich2010authdata-generic}: the generic
  framework for authenticated query answers; our holonomy ledger
  (\S 7, Sprint K) is a bundle-specialized authenticated data
  structure with Merkle root commitment and per-leaf
  $\mathrm{record\_hash}$ binding.
\end{itemize}

\subsection{The fiber-bundle data model}
\label{sec:background:bundle}

A bundle $B = (E, B, F, \pi, \Phi)$ stores records as sections of the
bundle. The base space $B$ is parameterized by the GIGI hash
$G: K_1 \times \cdots \times K_m \to \mathbb{Z}_2^{64}$ of the record's
primary keys. The fiber $F = F_1 \times \cdots \times F_k$ is the
Cartesian product of non-key field types. Write $\Gamma(E)$ for the
space of sections of $E \xrightarrow{\pi} B$. A record is a section
$\sigma \in \Gamma(E)$ whose pointwise image satisfies
$\sigma(p) = (p, v_1, \ldots, v_k) \in E$ for each base point $p$.
We will use $\sigma \in \Gamma(E)$ throughout when we mean ``a
section'' and $e \in E$ when we mean a point in the total space; on
encrypted writes the encryption map $\mathrm{Enc}_g$ acts on the
section by transforming each pointwise fiber value, hence
$\mathrm{Enc}_g: \Gamma(E) \to \Gamma(E)$.

\begin{definition}[Dispersion ratio $K$]
\label{def:dispersion}
For a numeric fiber field $F_i$ with finite-sample variance
$\mathrm{Var}_i$ and range $\mathrm{range}_i = \max_i - \min_i > 0$,
define the per-field \emph{dispersion ratio}
\[
K_i \;=\; \frac{\mathrm{Var}_i}{\mathrm{range}_i^2}\,.
\]
The bundle-level dispersion is the unweighted mean over numeric
fields, $K = \frac{1}{|\mathcal{N}|}\sum_{i \in \mathcal{N}} K_i$.
\end{definition}

\begin{remark}[Naming]
We sometimes call $K$ the bundle's ``scalar curvature'' following the
Davis~Manifold~\cite{davis2026manifold} naming convention, where
the dispersion ratio is identified with a discrete scalar curvature on
the section-graph manifold of the Davis substrate. The identification
is not assumed in this paper. Where $K$ appears in any claim or proof
here, only the definition above is used.
\end{remark}

\begin{definition}[Mean holonomy $\overline{\mathrm{Hol}}$]
\label{def:hol-mean}
Let $\Gamma_{\!B}$ be the discrete base graph of the bundle
(vertices = distinct base points, edges = field-index adjacencies
induced by the GIGI hash), and let $\beta_0, \beta_1$ be its $0$th
and $1$st Betti numbers (number of connected components; number
of independent cycles). The \emph{mean holonomy} is the
base-graph topological ratio
\[
\overline{\mathrm{Hol}}(B) \;=\; \frac{\beta_1}{\beta_0 + 1}.
\]
This is a base-only quantity: it depends purely on the base graph's
topology, not on fiber values. The motivation for naming it ``mean
holonomy'' is the proxy relationship to gauge-group holonomy:
non-trivial gauge holonomy is possible only when $\beta_1 > 0$
(cycles exist), and is increasingly likely as $\beta_1$ grows
relative to the number of components. The ratio is therefore an
upper-bound proxy for holonomic richness that is computable from
the base graph alone (no fiber decryption required).
\end{definition}

\begin{remark}[Gauge-group holonomy as a separate primitive]
A separate query, the GQL-level \texttt{HOLONOMY \ldots ON FIBER}
statement, computes the discrete Gauss-Bonnet angle-deficit
holonomy in a $2$-fiber $(f_0, f_1)$-plane around a chosen base
field. That operation reads stored fiber values; under v0.3.0 the
deficit it returned was gauge-dependent (per-field Aff with
$a_0 \neq a_1$ distorts angle calculations). \textbf{v0.3.1
normalizes centroids by their own min/range per axis before
computing angles}, which makes the deficit gauge-invariant
modulo $2\pi$ under any per-field Aff$(\mathbb{R})$ gauge --- the
fix is in
\texttt{src/bin/gigi\_stream.rs::compute\_fiber\_holonomy} and is
exercised by \texttt{tests/holonomy\_gauge\_invariance\_v0\_3.rs}
across a 20-gauge random sweep (including per-axis scale
mismatches, negative scales, and large translations). The
\texttt{HOLONOMY} statement therefore no longer requires
decryption: it can be evaluated on ciphertext and returns the
same deficit (mod $2\pi$) as it would on plaintext.

Throughout this paper, $\overline{\mathrm{Hol}}$ refers to the
base-graph topological ratio of
Definition~\ref{def:hol-mean} ($\beta_1 / (\beta_0 + 1)$), which is
the invariant-ring quantity reported by
\texttt{HolonomyAvg} in
\texttt{src/invariant.rs}~\cite{gigi2026code}. The \texttt{HOLONOMY
ON FIBER} GQL statement is a related but distinct primitive
(geometric Gauss-Bonnet on a fiber plane), not the
$\beta$-Betti-ratio that appears in the invariant tuple. The
detailed cases for $\overline{\mathrm{Hol}}$:
\begin{itemize}[nosep]
\item Tree-like base ($\beta_1 = 0$): $\overline{\mathrm{Hol}} = 0$.
\item One connected component with $k$ cycles:
  $\overline{\mathrm{Hol}} = k/2$.
\item Multiple components, no cycles: $\overline{\mathrm{Hol}} = 0$.
\end{itemize}
\end{remark}

\begin{definition}[Invariant tuple]
\label{def:invariant-tuple}
The bundle's six-component invariant fingerprint:
\[
\pi_{\mathrm{inv}}(B) \;=\;
\bigl( K,\; \lambda_1,\; \overline{\mathrm{Hol}},\; \tau,\; \beta_0,\; \beta_1 \bigr)
\]
where $K$ is the dispersion ratio (Definition~\ref{def:dispersion}),
$\lambda_1 := \lambda_1(L)$ is the smallest non-zero eigenvalue of
the combinatorial graph Laplacian $L = D - A$ on the field-index
bitmap graph $\mathcal{G}_{\mathrm{fi}}$ (the spectral gap;
equivalently the algebraic connectivity à la Fiedler),
$\overline{\mathrm{Hol}}$ is the mean holonomy
(Definition~\ref{def:hol-mean}), $\tau$ is the record count,
and $\beta_0, \beta_1$ are the $0$th and $1$st Betti numbers of
$\mathcal{G}_{\mathrm{fi}}$. (The v0.3.0 draft of this tuple included
$C = \tau/K$ instead of $\tau$ directly; see §3.5 for the rationale
for the v0.3.1 replacement.) This tuple is the load-bearing object
for Curvature-MAC (\S 5).
\end{definition}

\begin{remark}[What the Curvature-MAC tag does and does not protect]
\label{rmk:cmac-scope}
Curvature-MAC (\S 5) HMACs the canonical encoding of
$\pi_{\mathrm{inv}}(B)$ under a key derived from the gauge.
It detects edits to any of $(K, \lambda_1, \overline{\mathrm{Hol}},
\tau, \beta_0, \beta_1)$ exceeding the $10^{-10}$ quantization grain
--- i.e., \emph{any gauge-invariant content change}: record
insertion, deletion, or value change large enough to perturb a
component. It does \emph{not} on its own detect (a) ciphertext edits
that preserve all six invariants (a permutation of records, a
$+\varepsilon$ shift below the quantization grain) --- the
extended-leaf ledger of Sprint K (\S 7), which signs a per-record
\texttt{record\_hash}, closes this gap byte-level; (b) replay of an
old valid bundle under the same gauge --- the ledger's
append-only timestamp ordering closes this; (c) edits to
non-bundle data on the wire --- standard TLS / AEAD authentication
covers this layer. The pairing
(Curvature-MAC, ledger) is by design layered: invariant-level
attestation + per-record byte attestation, the way RFC 6962 layers
SCTs over per-certificate hashes.
\end{remark}

\subsection{Encryption modes and their mathematical structure}
\label{sec:background:groups}

\begin{definition}[Gauge action of an encryption mode]
\label{def:gauge-action}
An \emph{encryption mode} is a triple $(\mathcal{M}, \alpha, \mathcal{K})$
where $\mathcal{M}$ is a mathematical structure (a group, a quotient
map, or a randomized algorithm; see Table~\ref{tab:taxonomy} below),
$\alpha$ is the per-record action of $\mathcal{M}$ on a fiber tuple
$(v_1, \ldots, v_k) \in F$, and $\mathcal{K}$ is the key space
parameterizing $\mathcal{M}$. The base-space coordinates (keys) are
unchanged by $\alpha$.
\end{definition}

\begin{remark}
For the two deterministic-bijective modes (Affine and Isometric),
$\mathcal{M}$ is a genuine structure group $G \subset \mathrm{Aut}(F)$
acting linearly on $F$. For the three non-deterministic-bijective
modes (Opaque, Indexed, Probabilistic), $\mathcal{M}$ is a wider
mathematical object: Opaque uses a randomized AEAD, so the same
plaintext maps to a probability distribution over ciphertexts;
Indexed uses a deterministic but one-way PRF, so the action is a
quotient map rather than a bijection; Probabilistic combines a
bijective Aff($\mathbb{R}$) gauge with additive Gaussian noise.
Throughout the paper, we use ``structure group'' interchangeably
with ``encryption structure'' for the deterministic-bijective modes,
and specify the mathematical object explicitly when the mode is not
a group action.
\end{remark}

\begin{table}[h]
\centering
\small
\begin{tabular}{p{2cm}p{9cm}}
\toprule
\textbf{Mode} & \textbf{Mathematical structure $\mathcal{M}$} \\
\midrule
Affine & Structure group $G = \mathrm{Aff}(\mathbb{R})^k$ acting field-wise: $g_i(v) = a_i v + b_i$ with $a_i \neq 0$. \\
Isometric & Structure group $G = O(k)$ on a grouped $k$-tuple of numeric fields: $g(v) = O v + b$, $O^\top O = I$. \\
Opaque & Randomized AEAD on the byte encoding of a fiber field. Same plaintext maps to a distribution over ciphertexts (per-record nonce). Not a deterministic gauge; not invertible without the AEAD key. \\
Indexed & Deterministic PRF $F_k: F \to \{0,1\}^n$ keyed by $k$. The induced equivalence relation $v_1 \sim v_2 \Leftrightarrow F_k(v_1) = F_k(v_2)$ is the structure preserved on ciphertext; the encryption is a quotient map, not a bijection. \\
Probabilistic & Aff($\mathbb{R}$) gauge $(a, b)$ plus i.i.d.~additive noise: $g(v) = a v + b + \varepsilon$ with $\varepsilon \sim \mathcal{N}(0, \sigma^2)$. Bijective in expectation, randomized per-record. \\
\bottomrule
\end{tabular}
\caption{Encryption-mode taxonomy. Affine and Isometric carry a true
structure group; Opaque is randomized; Indexed is a deterministic
quotient; Probabilistic combines an Aff group action with additive
Gaussian noise.}
\label{tab:taxonomy}
\end{table}

The fundamental claim of \S\ref{sec:gauge-invariance}: the query class
computable on encrypted data depends on which sub-structure of
$\mathcal{M}$ is preserved by the encryption. For group-action modes,
this reduces to the structure group's invariant theory; for
randomized or quotient-map modes, it requires a separate
characterization per mode (which we develop in \S\S 4--9).


%% ==================================================================
\section{The gauge-equivariance theorem}
\label{sec:gauge-invariance}
%% ==================================================================

\subsection{Setup}
\label{sec:gauge-invariance:setup}

Let $B = (E, B, F, \pi, \Phi)$ be a fiber-bundle data store with
structure group $G \subset \mathrm{Aut}(F)$, and let $\Gamma(E)$
denote the space of sections (the bundle's record set; §2.3). Let
$g \in G$ be a secret gauge transformation. Let
$\mathrm{Enc}_g: \Gamma(E) \to \Gamma(E)$ be the geometric encryption
under $g$:
\[
\mathrm{Enc}_g(\sigma)(p) \;=\; \bigl(p,\, g(v_1, \ldots, v_k)\bigr)
\]
(base point $p$ unchanged; only fiber values transformed). Let
$f: \Gamma(E) \to \mathbb{R}$ be a query function on the bundle.

\begin{definition}[Gauge invariance]
\label{def:gauge-invariance}
$f$ is \emph{gauge-invariant under $G$} iff for every $g \in G$ and
every section $\sigma \in \Gamma(E)$:
\[
f(\mathrm{Enc}_g(\sigma)) \;=\; f(\sigma).
\]
\end{definition}

\subsection{The framework: equivariance and ciphertext computability}
\label{sec:gauge-invariance:thm}

The reviewer-friendly statement of the framework is not an
invariance characterization but an \emph{equivariance}
characterization. Many of the queries an analyst actually wants ---
$\textsc{sum}$, $\textsc{avg}$, $\textsc{min}$, $\textsc{max}$,
$\textsc{var}$, $\textsc{stddev}$, $\textsc{range}$ over affine
ciphertext --- are not gauge-invariant (the encrypted sum is
\emph{not} the plaintext sum). They are
\emph{$G$-equivariant} via a closed-form representation $\rho$, and
the client recovers the plaintext answer by applying $\rho_g^{-1}$
once at $O(1)$ cost. The invariant case is the special instance
$\rho \equiv \mathrm{id}$. We state the unified theorem first; the
invariant ring of \S\,\ref{sec:gauge-invariance:ring} and the
client-inversion helpers of \S\,\ref{sec:gauge-invariance:examples}
are both corollaries.

\begin{definition}[Equivariant query]
\label{def:equivariant-query}
Let $G$ act on $\Gamma(E)$ via $\mathrm{Enc}_g$
(Definition~\ref{def:gauge-action}), let
$f: \Gamma(E) \to \mathbb{R}^m$ be a query function, and let
$\rho: G \to \mathrm{Aut}(\mathbb{R}^m)$ be a group representation.
$f$ is \emph{$\rho$-equivariant under $G$} iff for every $g \in G$
and every $\sigma \in \Gamma(E)$:
\[
f(\mathrm{Enc}_g(\sigma)) \;=\; \rho_g \cdot f(\sigma).
\]
When $\rho_g \equiv \mathrm{id}_{\mathbb{R}^m}$, $f$ is
\emph{gauge-invariant} (Definition~\ref{def:gauge-invariance}, the
$m=1$, $\rho = \mathrm{triv}$ instance).
\end{definition}

\begin{theorem}[Equivariant ciphertext-computability]
\label{thm:equivariance}
Let $G$ act on $\Gamma(E)$ via $\mathrm{Enc}_g$ with
$\mathrm{Enc}_e = \mathrm{id}_{\Gamma(E)}$, and let
$\rho: G \to \mathrm{Aut}(\mathbb{R}^m)$ be a representation. For
$f: \Gamma(E) \to \mathbb{R}^m$ the following are equivalent:
\begin{enumerate}[label=(\alph*),nosep]
\item \textbf{(Equivariance.)} $f$ is $\rho$-equivariant:
  $f(\mathrm{Enc}_g(\sigma)) = \rho_g(f(\sigma))$ for all $g, \sigma$.
\item \textbf{(Ciphertext-side computability up to $\rho$.)} There
  exists a deterministic algorithm $A_\rho: \Gamma(E) \to \mathbb{R}^m$
  such that
  $A_\rho(\mathrm{Enc}_g(\sigma)) = \rho_g(f(\sigma))$ for all
  $g \in G, \sigma \in \Gamma(E)$.
\item \textbf{(Client-side recovery with $O(\mathrm{poly}(m))$ post-processing.)}
  There exists $A_\rho$ as in (b) such that the client, given the
  secret gauge $g$, recovers the plaintext answer via
  \[
  f(\sigma) \;=\; \rho_g^{-1}\!\bigl(A_\rho(\mathrm{Enc}_g(\sigma))\bigr),
  \]
  at cost $O(\mathrm{poly}(m))$ \emph{independent of $|E|$}.
\end{enumerate}
The invariant special case $\rho \equiv \mathrm{id}$ collapses
(b)/(c) to $A(\mathrm{Enc}_g(\sigma)) = f(\sigma)$ --- the
no-post-processing regime --- and recovers the v0.3.0 framing of
\S\,3.2 as a corollary (Corollary~\ref{cor:invariant-special-case}).
\end{theorem}

\begin{proof}
\emph{(a) $\Rightarrow$ (b).} Take $A_\rho = f$. For any $\sigma, g$:
$A_\rho(\mathrm{Enc}_g(\sigma)) = f(\mathrm{Enc}_g(\sigma))
= \rho_g(f(\sigma))$ by (a). So $A_\rho = f$ witnesses (b).

\emph{(b) $\Rightarrow$ (c).} Since $\rho_g \in \mathrm{Aut}(\mathbb{R}^m)$
is invertible, $\rho_g^{-1}$ is well-defined; the client computes it
from $g$ in $O(m^3)$ (matrix inversion when $m > 1$; scalar
reciprocation when $m = 1$) and applies it to the server's output.
The cost is in $m$, the answer dimension, not $|E|$.

\emph{(c) $\Rightarrow$ (a).} Let $A_\rho$ witness (c). Fix any
$\sigma$ and $g$, and let $\sigma' := \mathrm{Enc}_g(\sigma)$. Apply
$A_\rho$ to $\sigma'$ along two paths:
\begin{itemize}[nosep]
\item Path 1 (true gauge $g$): $\sigma' = \mathrm{Enc}_g(\sigma)$, so
  by (c): $A_\rho(\sigma') = \rho_g(f(\sigma))$.
\item Path 2 (identity gauge $e$): $\sigma' = \mathrm{Enc}_e(\sigma')$
  since $\mathrm{Enc}_e = \mathrm{id}_E$, and by (c) applied to
  $\sigma'$ and gauge $e$: $A_\rho(\sigma') = \rho_e(f(\sigma'))
  = f(\sigma')$ because $\rho$ is a homomorphism so $\rho_e
  = \mathrm{id}_{\mathbb{R}^m}$.
\end{itemize}
Equating: $\rho_g(f(\sigma)) = f(\sigma')
= f(\mathrm{Enc}_g(\sigma))$. Since $g, \sigma$ were arbitrary,
$f$ is $\rho$-equivariant.
\end{proof}

\begin{corollary}[Invariant special case]
\label{cor:invariant-special-case}
With $\rho \equiv \mathrm{id}$ and $m = 1$,
Theorem~\ref{thm:equivariance} reduces to: $f$ is $G$-invariant iff
there exists $A: E \to \mathbb{R}$ with
$A(\mathrm{Enc}_g(\sigma)) = f(\sigma)$ for all $g, \sigma$. This is
the ``invariance characterizes ciphertext computability'' framing of
v0.3.0; it is now a corollary of the equivariant theorem.
\end{corollary}

\begin{remark}[Why equivariance is the right framework]
The earlier v0.3.0 framing as ``invariance characterizes
ciphertext computability'' systematically undercounted the
ciphertext-computable query class. $\textsc{sum}$ on an
affine-encrypted column has $\rho_g(s) = a \cdot s + n \cdot b$ ---
a non-trivial $\rho$, so SUM is \emph{not} invariant but \emph{is}
ciphertext-computable: the encrypted sum, post-processed by
$\rho_g^{-1}(\hat s) = (\hat s - n b)/a$ at the client, equals the
plaintext sum exactly. The same closed-form $\rho_g^{-1}$ recovers
AVG, MIN, MAX, VAR, STDDEV, RANGE
(\S\,\ref{sec:gauge-invariance:examples}). Calling these queries
``not computable from the ciphertext'' would have been a strict
factual error; the equivariance lens repairs it.
\end{remark}

\begin{remark}[Closed-form vs.\ general decryption]
The post-processing in (c) is $\rho_g^{-1}$, a closed-form
algebraic operation derived from $g$ and the query's
representation $\rho$. It is \emph{not} full decryption of the
bundle, and it is \emph{not} computation on plaintext; the client
applies a single $O(\mathrm{poly}(m))$ matrix (typically $m = 1$,
so a single scalar division). The server's compute remains over
ciphertext at native speed, on a fully encrypted bundle.
\end{remark}

\begin{remark}[Equivariant ring vs.\ invariant ring]
For each structure group $G$ and representation $\rho$,
Theorem~\ref{thm:equivariance} identifies the
\emph{$\rho$-equivariant subspace}
$\mathcal{E}_{G,\rho} \subset C^\infty(E, \mathbb{R}^m)$. The
invariant ring $\mathcal{I}_G = \mathcal{E}_{G,\mathrm{id}}$ of
\S\,\ref{sec:gauge-invariance:ring} is the trivial-representation
slice. The non-trivial slices give the analytical SQL surface
(SUM/AVG/etc.) developed in \S\,\ref{sec:gauge-invariance:examples}.
This is the database-runtime analogue of equivariant
representation theory: the ciphertext-computable algebra
\emph{decomposes by irreducible representation of $G$} into a
direct sum of equivariant subspaces, exactly one of which
(the trivial summand) is the invariant ring.
\end{remark}

\begin{theorem}[Affine-mode invariant ring contains the bundle's invariant tuple]
\label{thm:affine-invariants}
Let $G = \mathrm{Aff}(\mathbb{R})^k$ act field-wise on the fiber
$F = \mathbb{R}^k$. The six components of $\pi_{\mathrm{inv}}(B)$
(Definition~\ref{def:invariant-tuple}) lie in
$\mathcal{I}_G$, i.e.~each is fixed under any per-field affine
re-encryption with $a_i \neq 0$.
\end{theorem}

\begin{proof}
We check each component:
\begin{itemize}[nosep]
\item $K = \frac{1}{|\mathcal{N}|}\sum_i \mathrm{Var}_i / \mathrm{range}_i^2$.
  Under $g_i(v) = a_i v + b_i$: $\mathrm{Var}_i(g_i(v)) = a_i^2 \mathrm{Var}_i(v)$,
  $\mathrm{range}_i(g_i(v)) = |a_i| \cdot \mathrm{range}_i(v)$, so the
  ratio is invariant. The mean over fields is invariant.
\item $\lambda_1$: the spectral gap is computed from the field-index
  bitmap topology (Definition~\ref{def:invariant-tuple}), which is a
  binary indicator of which records share which values. Affine
  per-field encryption is injective on each field, so the
  value-equivalence classes are preserved, hence the bitmap topology
  is preserved, hence $\lambda_1$ is bit-identical.
\item $\overline{\mathrm{Hol}} = \beta_1 / (\beta_0 + 1)$
  (Definition~\ref{def:hol-mean}) depends purely on the base graph's
  Betti numbers, which are preserved under affine per-field encryption
  (as in the $\lambda_1$ case). $\overline{\mathrm{Hol}}$ is therefore
  preserved bit-identically.
\item $\tau$: record count is unchanged by any deterministic
  encryption.
\item $\beta_0$, $\beta_1$: graph-topology invariants on the
  field-index bitmap; preserved under value-injective gauges as above.
\end{itemize}
All six components are therefore fixed under $G$.
\end{proof}

\begin{remark}[On generating sets]
Theorem~\ref{thm:affine-invariants} states containment in
$\mathcal{I}_{\mathrm{Aff}}$, not generation. The question of
whether these six scalars generate the full invariant ring is
mode-dependent and outside this paper's scope; we make claims only
about containment and about the gauge-rotation invariance of
$\pi_{\mathrm{inv}}$ under the listed modes.
\end{remark}

\subsection{Examples per mode}
\label{sec:gauge-invariance:examples}

\paragraph{Affine ($G = \mathrm{Aff}(\mathbb{R})^k$).}
\emph{Invariant slice} ($\rho_g = \mathrm{id}$): the six
components of $\pi_{\mathrm{inv}}$
(Theorem~\ref{thm:affine-invariants}); plus order and equality
predicates --- the latter via the affine path
$\textsc{where}\ v > c \mapsto \textsc{where}\ g(v) > g(c)$
(transform the literal, not the data; order is preserved iff $a > 0$,
equality is always bijective).
\smallskip

\noindent
\emph{Equivariant slice} (non-trivial $\rho_g$): the analytical-SQL
aggregates SUM, AVG, MIN, MAX, VAR, STDDEV, RANGE are computable on
ciphertext at native server speed, then inverted at the client in
$O(1)$ via $\rho_g^{-1}$. For a single column under
$g(v) = a v + b$ over $n$ records:
\[
\begin{aligned}
&\rho_{\textsc{sum}}(s) = a s + n b,
\quad
&\rho_{\textsc{avg}}(\mu) = a \mu + b,\\
&\rho_{\textsc{min}}(m) = a m + b \;(a>0),
\quad
&\rho_{\textsc{max}}(M) = a M + b \;(a>0),\\
&\rho_{\textsc{var}}(V) = a^2 V,
\quad
&\rho_{\textsc{stddev}}(s) = |a| s,\\
&\rho_{\textsc{range}}(R) = |a| R.
\end{aligned}
\]
Each $\rho_g$ is a closed-form invertible map in
$\mathrm{Aut}(\mathbb{R})$; the client inversion
$\rho_g^{-1}$ is one division and one subtraction. This is the
equivariant ring under
$\rho: \mathrm{Aff}(\mathbb{R}) \to \mathrm{Aut}(\mathbb{R})$ for
each of the seven aggregates; the reference implementation is
\texttt{src/aggregate\_helpers.rs} ($16$ unit tests, $12$ exactness
+ $4$ bias-refusal under Probabilistic-mode noise).

\paragraph{Opaque (randomized AEAD).}
Functions of the encrypted fiber values are not invariant
(only constants), because AEAD output is a randomized function of
plaintext + nonce. Functions of \emph{the base alone} --- record
count $\tau$, Betti numbers $\beta_0$, $\beta_1$, schema structure
--- are trivially invariant because the encryption does not touch
the base space (Definition~\ref{def:gauge-action}).
Computability class on the fiber: nothing without decryption.

\paragraph{Indexed (deterministic PRF).}
$v_1 = v_2 \Rightarrow \mathrm{PRF}_k(v_1) = \mathrm{PRF}_k(v_2)$
deterministically; the converse holds with overwhelming probability
$1 - 2^{-128}$ for the AES-CMAC instantiation, the standard
PRF-collision floor at the $128$-bit security level (so spurious
collisions on the encrypted side are negligible and the verifier
treats them as an honest content edit). Computability: equality
and equality-based queries (\textsc{group by}, \textsc{join} on
equality, bitmap-index lookup). Functions of the base also remain
trivially invariant. The structure group is degenerate (a
one-way quotient, not a group action), so
Theorem~\ref{thm:equivariance} applies in the form ``the
equivariance class is the equality predicate $v_1 = v_2$,
realized through the PRF rather than through a closed-form
$\rho_g$.''

\paragraph{Probabilistic ($g(v) = av + b + \varepsilon$ with $\varepsilon \sim \mathcal{N}(0, \sigma^2)$ per encryption).}
Probabilistic mode breaks exact equivariance: the per-encryption
noise $\varepsilon$ is fresh on every ciphertext, so
$\mathrm{Enc}_g$ is not deterministic and no $\rho_g$ literally
satisfies
$f(\mathrm{Enc}_g(\sigma)) = \rho_g(f(\sigma))$ pointwise.
Theorem~\ref{thm:equivariance} therefore does not apply
directly --- it is replaced by an \emph{$\epsilon$-equivariance}
characterization in the concentration sense:
\[
\Pr\Bigl[\,\bigl|\,f(\mathrm{Enc}_g(\sigma))
- \rho_g(f(\sigma))\,\bigr| \;\geq\; \epsilon\,\Bigr]
\;\leq\; \delta(\epsilon, n, \sigma),
\]
where $n$ is the records-in-aggregate count and
$\delta(\epsilon, n, \sigma)$ is a concentration tail. For the
linear aggregates SUM / AVG with $g(v) = av + b$ and $n$
independent Gaussian noise draws:
\begin{itemize}[nosep]
\item \textbf{SUM.}\quad
  $\widehat{S}_n - (a S_n + n b) = \sum_{i=1}^n \varepsilon_i \sim
  \mathcal{N}(0, n \sigma^2)$. Hoeffding / Chernoff yields
  \[
  \Pr\!\left[\,\bigl|\widehat{S}_n - (aS_n + nb)\bigr| \geq \epsilon\,\right]
  \;\leq\; 2 \exp\!\left(-\frac{\epsilon^2}{2 n \sigma^2}\right),
  \]
  so for $\epsilon = c \sigma \sqrt{n}$ the failure probability is
  $\leq 2 e^{-c^2/2}$ --- standard sub-Gaussian concentration.
\item \textbf{AVG.}\quad The mean error
  $\widehat{\mu}_n - (a \mu_n + b) \sim \mathcal{N}(0, \sigma^2/n)$
  contracts at rate $1/\sqrt{n}$:
  $\Pr[|\widehat{\mu}_n - (a\mu_n + b)| \geq \epsilon]
  \leq 2 \exp(-n \epsilon^2 / (2\sigma^2))$.
\item \textbf{VAR.}\quad Bernstein's inequality applies to the
  bounded centered moments, giving a one-sided tail
  $\Pr[\widehat{V}_n - (a^2 V_n + \sigma^2) \geq \epsilon]
  \leq \exp\!\bigl(-n\epsilon^2/(2(2\sigma^4 + 2\sigma^2 M^2 + M \epsilon/3))\bigr)$
  where $M$ is an a-priori bound on $|av+b|$ in the bundle.
\end{itemize}
The result is an $(\epsilon,\delta)$-equivariant ring under $\rho$
covering \emph{linear and quadratic aggregates}
(COUNT, SUM, AVG, VAR, STDDEV) with closed-form correction
$\rho_g^{-1}$ on the noisy estimator: the
client recovers the plaintext aggregate to within an additive
$O(\sigma/\sqrt{n})$ error with high probability.
\emph{Order statistics (MIN, MAX, RANGE) are excluded}: the
extreme-value functional does not commute with additive Gaussian
noise, and the naïve closed-form $\rho_g^{-1}$ overshoots in
expectation by $\Theta(\sigma\sqrt{2\log n}/|a|)$, a bias that
does \emph{not} contract as $n \to \infty$. \S\,\ref{sec:intro:not}
documents the resulting deployment caveat and the three
operational mitigations (use VAR-derived spread, switch to Affine
$\sigma = 0$, or escalate the column to FHE); the shipped helpers
refuse \texttt{decrypt\_min}/\texttt{\_max}/\texttt{\_range} under
$\sigma > 0$ with a typed \texttt{BiasedUnderProbabilisticNoise}
error so the bias cannot leak into application code unnoticed.
For $\sigma$-bucket equality predicates, the Davis Identity
$v_1 \approx_\sigma v_2 \Leftrightarrow
\mathrm{bucket}_\sigma(v_1) = \mathrm{bucket}_\sigma(v_2)$
gives a deterministic equivalence relation on the noised
ciphertext space, so equality-modulo-$\sigma$ remains an
exact invariant predicate.

\paragraph{Isometric ($G = E(k) = O(k) \ltimes \mathbb{R}^k$ on a grouped $k$-tuple).}
We take the full Euclidean group $E(k) = O(k) \ltimes \mathbb{R}^k$
(rotations / reflections composed with translations) as the
structure group on a grouped $k$-tuple; the orthogonal-only case
$G = O(k)$ is the translation-fixed subgroup and is recovered by
setting the translation to $0$. Pairwise Euclidean distance
$\|u - v\|$ is preserved exactly under $E(k)$ --- both rotation
and translation preserve distances --- hence any distance-based
query (k-NN, clustering, distance threshold) is invariant.

Note that the dispersion ratio $K_i$
(Definition~\ref{def:dispersion}) is \emph{per-field} and is in
general \emph{not} preserved by $E(k)$ --- a rotation in $O(2)$ on
fields $(x, y)$ mixes the per-field variances and ranges, and a
translation shifts the per-field mins/maxes. The
$E(k)$-invariant generalization of dispersion is the
trace-form / diameter quotient
\[
K_{\mathrm{group}}(X) \;=\; \frac{\mathrm{tr}\,\mathrm{Cov}(X)}
                              {\mathrm{diam}^2(X)},
\qquad
\mathrm{diam}^2(X) \;:=\; \max_{p,q \in X}\, \|p - q\|^2,
\]
where $X = \{v^{(1)}, \ldots, v^{(n)}\} \subset \mathbb{R}^k$ is
the set of record-valued grouped vectors. Both numerator and
denominator are $E(k)$-invariant: the trace of the covariance
matrix is the sum of eigenvalues, which is conjugation-invariant
under $O(k)$ and unchanged by mean-shifts (translations); the
squared diameter is a max over pairwise Euclidean distances and is
preserved exactly under $E(k)$ by definition.

The earlier v0.3.0 draft used a per-coordinate range
$(\max_i v - \min_i v)^2$ in the denominator; this is \emph{not}
$O(k)$-invariant (the coordinate-wise extrema change under
rotation) and was a bug. The squared-diameter denominator
$\mathrm{diam}^2(X)$ above is the corrected, genuinely
$E(k)$-invariant form.

The Curvature-MAC construction
(\S\,5) uses the per-field $K$ for fields
\emph{outside} an isometric group and $K_{\mathrm{group}}$ for
the group's grouped components; see \S\,5
for the precise mode-conditional definition. Mean holonomy
$\overline{\mathrm{Hol}}$ (Definition~\ref{def:hol-mean}) depends
on the field-index topology rather than the value coordinates and
is preserved by $E(k)$ (the topology is unaffected by isometric
re-coordinatization of values).

\needspace{6\baselineskip}
\paragraph{Per-mode leakage profile.}
The honest summary of \emph{what each mode reveals to a
ciphertext-only adversary} is the table below. Leakage is graded by
the standard structured-encryption taxonomy of
Chase--Kamara~\cite{chase2010structured}: \emph{access pattern}
$\mathcal{L}_{\mathrm{acc}}$ (which records were touched by which
query), \emph{equality pattern} $\mathcal{L}_{\mathrm{eq}}$ (which
pairs of records share a value on a queried field), \emph{order
pattern} $\mathcal{L}_{\mathrm{ord}}$ (the value-ordering on a
queried field), and \emph{distance pattern}
$\mathcal{L}_{\mathrm{dist}}$ (pairwise Euclidean distances within
an isometrically-grouped tuple).
\begin{center}\footnotesize
\renewcommand{\arraystretch}{1.2}
\setlength{\tabcolsep}{4pt}
\begin{tabularx}{\linewidth}{|l|c|c|c|c|X|}
\hline
\textbf{Mode} & \textbf{Eq.} & \textbf{Ord.} & \textbf{Dist.} & \textbf{Aggregate} & \textbf{Adversary class} \\
\hline
Opaque (AEAD)              & $\bot$ & $\bot$ & $\bot$ & $\bot$ & IND-CPA / IND-CCA \\
Indexed (det.\ PRF)        & $\mathcal{L}_{\mathrm{eq}}$ & $\bot$ & $\bot$ & $\bot$ & PRF-secure, $2^{-128}$ collision floor \\
Affine (Aff$(\mathbb{R})$) & $\mathcal{L}_{\mathrm{eq}}$ & $\mathcal{L}_{\mathrm{ord}}$ & $\bot$ & \textbf{equiv.}~via $\rho_g$ & IND-OCPA (single column) \\
Probabilistic              & $\approx_\sigma$ & $\bot$ & $\bot$ & $(\epsilon,\delta)$-equiv.~via $\rho_g$ & noise-bucket OCPA ($\sigma$) \\
Isometric ($E(k)$)         & $\bot$ & $\bot$ & $\mathcal{L}_{\mathrm{dist}}$ & dist-invariant & distance-only (no per-coord) \\
\hline
\end{tabularx}
\end{center}
Three things to read out of this table:
\begin{itemize}[nosep]
\item Each mode's leakage profile is the \emph{minimal} information
  the gauge-equivariance structure forces the server to learn.
  Anything beyond this is a leakage bug, not a leakage feature.
\item \textbf{Equivariant computability is a separate axis from
  leakage}: the Affine and Probabilistic rows compute SUM / AVG /
  MIN / MAX / VAR / STDDEV / RANGE via the equivariant representation
  $\rho_g$ \emph{without expanding the access / order / equality
  leakage profile} --- the client's $\rho_g^{-1}$ is purely
  post-processing of the server's aggregate output.
\item All five modes share access-pattern leakage
  $\mathcal{L}_{\mathrm{acc}}$ at the bundle layer (which records
  are read by which query); the standard
  oblivious-RAM / private information retrieval mitigations apply
  uniformly and are orthogonal to the gauge-invariance
  construction.
\end{itemize}

The forthcoming \S\,4 develops the per-mode security proofs
(reductions to IND-CPA, PRF, BDH, MLWE) against the leakage
profiles above; the present subsection's purpose is to make the
profile explicit before the proofs land.

\subsection{The invariant ring}
\label{sec:gauge-invariance:ring}

For each deterministic-bijective mode with structure group $G$, the
$G$-invariant ring is the subalgebra of real-valued bundle functions
fixed by $G$:
\[
\mathcal{I}_G \;=\; \bigl\{\, f \in \mathrm{Map}(E, \mathbb{R})
                          \;\big|\; f \circ \mathrm{Enc}_g = f
                          \text{ for all } g \in G \,\bigr\}.
\]
$\mathcal{I}_G$ is closed under pointwise $+$, $\times$, and
multiplication by scalars in $\mathbb{R}$. Theorem~\ref{thm:affine-invariants}
shows that for $G = \mathrm{Aff}(\mathbb{R})^k$,
$\mathcal{I}_{\mathrm{Aff}}$ contains the six components of
$\pi_{\mathrm{inv}}$ (Definition~\ref{def:invariant-tuple}). Further
$\mathcal{I}_{\mathrm{Aff}}$ contains every polynomial combination of
those components; whether they generate the full invariant ring is
left open here (see Remark on generating sets following
Theorem~\ref{thm:affine-invariants}). The Sprint H
\verb|PROJECT INVARIANT(...)| query language is the parser surface
that admits polynomial-combination expressions over $\pi_{\mathrm{inv}}$;
it is by construction a sublanguage of $\mathcal{I}_{\mathrm{Aff}}$.
This pattern --- restricting the query grammar to admit only
operations on a committed invariant subspace --- is the
database-runtime analogue of the privacy-preserving-blueprints
proof discipline~\cite{kohlweiss2023blueprints}, in which the
admissible statements about a committed value are restricted to
the policy-satisfying subspace.

\subsection{Floating-point reality}
\label{sec:gauge-invariance:fp}

Theorem~\ref{thm:equivariance} (and its invariant
specialization Corollary~\ref{cor:invariant-special-case}) hold in
exact arithmetic. In floating-point implementation, gauge
invariance/equivariance holds only \emph{up to} the precision of
the arithmetic. Specifically, re-encrypting a
bundle under a new gauge $g'$ (the operation Sprint G's
\verb|ROTATE_KEY| performs) applies one affine round-trip
$v \mapsto g(v) \mapsto g^{-1}(g(v)) \mapsto g'(g^{-1}(g(v)))$ to each
field, introducing $\sim 10^{-13}$ relative error per traversal.

The effect on the components of the invariant tuple
$\pi_{\mathrm{inv}}(B)$ (Definition~\ref{def:invariant-tuple}):
\begin{itemize}[nosep]
\item $\beta_0$, $\beta_1$, $\tau$ are integers (Betti numbers and
  record count) --- exact, zero drift.
\item $\lambda_1$ is computed from the field-index bitmap topology by
  integer-arithmetic eigenvalue routines on the $0$/$1$ adjacency ---
  exact under any value-preserving gauge transformation.
\item $K = \mathrm{Var}/\mathrm{range}^2$ drifts by $\sim 10^{-11}$
  relative through the Welford-streaming variance computation across
  the rotated record set --- the only f64 component with any drift.
\item $\overline{\mathrm{Hol}} = \beta_1 / (\beta_0 + 1)$ is a ratio
  of integer Betti numbers and is therefore \emph{exact}: zero
  drift under any gauge that preserves the base-graph topology.
\end{itemize}

\paragraph{Why $\tau$ replaces capacity in $\pi_{\mathrm{inv}}$.}
An earlier draft (v0.3.0) included $C = \tau/K$ explicitly in the
invariant tuple. This amplified $K$'s relative drift ($\sim 10^{-11}$)
into capacity's absolute drift by a factor of $\tau$: for a 40-record
bundle, the observed drift on $C$ was $\sim 4 \times 10^{-9}$ ---
$10^4 \times$ worse than $K$'s drift. Because $C$ is fully recoverable
on demand from $K$ and $\tau$, storing it in the tuple input adds no
information beyond what already binds the tag: $K$ is signed at full
precision, $\tau$ is a $u64$ with zero drift, and $C = \tau/K$ is
trivially recomputable by any caller. v0.3.1 therefore replaces the
$C$ slot with a $\tau$ slot in the canonical-bytes encoding. The
information content is identical; the drift amplification is gone.

\begin{remark}[Reading numbers across drafts]
The abstract's reference to ``$\sim 10^{-9}$ noise on the Davis
capacity'' is the v0.3.0 $C$-drift figure --- the resolved
amplification problem. The v0.3.1 worst remaining f64 component
drift is $K$'s $\sim 10^{-13}$ absolute (i.e.\ $\sim 10^{-11}$
relative), four orders of magnitude below the v0.3.0 number.
\end{remark}

\paragraph{Quantization grain.}
With $\overline{\mathrm{Hol}}$ now a Betti-ratio (integer-exact) and
$\lambda_1$ computed by integer-arithmetic eigenvalue routines on a
$0$/$1$ bitmap, the only f64 component with any drift is $K$. We
quantize the f64 slots ($K$, $\lambda_1$, $\overline{\mathrm{Hol}}$)
uniformly to $10$ decimal digits before HMAC:
\[
\mathrm{quantize}(x) \;=\; \mathrm{round}(x \cdot 10^{10}) \in \mathbb{Z}.
\]
The measured worst-case drift on any f64 component is now $\sim 10^{-13}$
absolute (on $K$), $3$ orders of magnitude below the $10^{-10}$
quantization grain. The framework's gauge-invariance
property (Theorem~\ref{thm:equivariance} +
Corollary~\ref{cor:invariant-special-case} +
Theorem~\ref{thm:affine-invariants}) therefore holds bit-identically
through the integrity tag in practice. The quantization grain is
chosen $3$ orders of magnitude looser than the measured drift so
that legitimate floating-point noise rounds away \emph{below} the
HMAC input; any genuine drift in the invariants reaching the
$10^{-10}$ digit --- $4$ orders of magnitude tighter sensitivity
than v0.3.0's $10^{-6}$ grain --- is preserved through the tag and
detected as a content edit by the verifier.

\paragraph{Reference implementation.}
The fix and measured drift numbers are reproducible from the public
GIGI repository. Three artifacts:
\begin{itemize}[nosep,leftmargin=1.5em]
\item \texttt{src/integrity.rs} --- the \texttt{InvariantTuple}
  struct and \texttt{quantize\_f64} helper.
\item \texttt{tests/composition\_v0\_3.rs} --- end-to-end roundtrip
  test\\
  \texttt{test\_composition\_\allowbreak{}integrity\_tag\_\allowbreak{}invariant\_under\_\allowbreak{}gauge\_rotation}\\
  on a $40$-record encrypted bundle.
\item \texttt{theory/encryption/validation/}\\
  \texttt{validation\_tests\_v0\_3.py} --- independent Python math
  oracle.
\end{itemize}

\subsection{Implications for the rest of the paper}
\label{sec:gauge-invariance:implications}

\S 4 catalogs the five modes with their mathematical structure. \S 5
builds Curvature-MAC: an HMAC on the invariant tuple,
gauge-rotation-invariant in floating-point thanks to the $10$-dp
quantization. \S 6 builds two delegation modes side-by-side:
Aff$(\mathbb{R})$ closure (trusted-delegatee, $O(1)$ algebraic
recovery on collusion; intended for honest-delegatee deployments
where Bob's key is independently protected) and pairing-based
key-delegation on BLS12-381 (collusion-resistance under DLP on
$G_2$; suitable when the delegatee or Bob's key may be
compromised). \S\S 7--9 build three further constructions on top.

The thread tying them together is
Theorem~\ref{thm:equivariance}: each construction either
\emph{preserves} a gauge invariant under some operation
(Curvature-MAC, capability delegation, ledger telescope), exploits
a non-trivial \emph{equivariance} for ciphertext-side computation
with closed-form client recovery (analytical SQL helpers
\S\,\ref{sec:gauge-invariance:examples}), or \emph{advances the
gauge} in a controlled way (RG-flow ratchet). The taxonomy of
what's possible flows directly from the choice of structure group
and the representation $\rho$ each query realises.


%% ==================================================================
%% Forward references and reproducibility (placeholder block for future sections)
%% ==================================================================

\section*{Forthcoming sections}

\S 4 through \S 14 will be drafted in subsequent revisions:

\begin{itemize}[nosep]
\item \S 4 --- The mode taxonomy (one subsection per mode).
\item \S 5 --- Curvature-MAC (Sprint I).
\item \S 6 --- $\mathrm{Aff}(\mathbb{R})$ capability delegation
  (Sprint J), including \textbf{Limitation 4.7.1 (collusion path
  documented explicitly)}.
\item \S 7 --- Holonomy ledger (Sprint K).
\item \S 8 --- \v{C}ech threshold sharing (Sprint L).
\item \S 9 --- Continuous RG-flow ratchet (Sprint M).
\item \S 10 --- Empirical evaluation (981+ Rust tests + 66 Python
  oracles).
\item \S 11 --- Deferred constructions (geodesic-ball ABE,
  spectral-signature ZKP, lattice-fiber PQ gauge).
\item \S 12 --- Related work.
\item \S 13 --- Limits.
\item \S 14 --- Conclusion.
\end{itemize}

The math primitives backing each forthcoming section are already
implemented in the public GIGI repository at
\url{https://github.com/nurdymuny/gigi/}; the Python math-validation
suite at \texttt{theory/encryption/validation/validation\_tests\_v0\_3.py}
captures $66$ independent oracle tests for the constructions in
\S\S 5--9.


%% ==================================================================
\section*{Acknowledgments}
%% ==================================================================

Solo-authored. Drafting and editing of this manuscript were assisted by
Claude (Anthropic); the mathematical positions, design choices, framing
decisions, and acceptance or rejection of empirical results are the
author's. The GIGI runtime referenced throughout is maintained at
Davis Geometric; the test surface backing every claim in this paper is
reproducible from the public \texttt{gigi} repository alone.


%% ==================================================================
\begin{thebibliography}{99}

\bibitem{davis2026sheafcomposition}
B.~R.\ Davis,
\emph{Sheaf Composition: The Geometry of Creativity, Implemented},
Zenodo, 2026, doi:10.5281/zenodo.20185331.

\bibitem{davis2026manifold}
B.~R.\ Davis,
\emph{The Davis Manifold: A Riemannian Substrate for Residue-Bounded
Computation},
Davis Geometric, 2026. Zenodo DOI to be assigned.

\bibitem{davis2026kahlersubstrate}
B.~R.\ Davis,
\emph{The K\"ahler Substrate: A Multi-Layer Geometric Engine for the
Davis Framework, Validated by Sheaf-Composition Runtime Observation},
Davis Geometric, 2026. Zenodo DOI to be assigned.

\bibitem{gigi2026code}
B.~R.\ Davis,
\emph{GIGI: Geometric Intrinsic Global Index} (source code),
Davis Geometric, 2026.
Available at \url{https://github.com/nurdymuny/gigi/}.

\bibitem{gentry2009fhe}
C.~Gentry,
Fully homomorphic encryption using ideal lattices,
in \emph{Proc.\ STOC 2009}, pp.\ 169--178.

\bibitem{bgv2012}
Z.~Brakerski, C.~Gentry, V.~Vaikuntanathan,
(Leveled) Fully homomorphic encryption without bootstrapping,
in \emph{Proc.\ ITCS 2012}, pp.\ 309--325.

\bibitem{ckks2017}
J.~H.\ Cheon, A.~Kim, M.~Kim, Y.~Song,
Homomorphic encryption for arithmetic of approximate numbers,
in \emph{Proc.\ ASIACRYPT 2017}, pp.\ 409--437.

\bibitem{boldyreva2009ope}
A.~Boldyreva, N.~Chenette, Y.~Lee, A.~O'Neill,
Order-preserving symmetric encryption,
in \emph{Proc.\ EUROCRYPT 2009}, pp.\ 224--241.

\bibitem{popa2011cryptdb}
R.~A.\ Popa, C.~M.~S.\ Redfield, N.~Zeldovich, H.~Balakrishnan,
CryptDB: protecting confidentiality with encrypted query processing,
in \emph{Proc.\ SOSP 2011}, pp.\ 85--100.

\bibitem{chase2010structured}
M.~Chase, S.~Kamara,
Structured encryption and controlled disclosure,
in \emph{Proc.\ ASIACRYPT 2010}, pp.\ 577--594.

\bibitem{curtmola2006sse}
R.~Curtmola, J.~Garay, S.~Kamara, R.~Ostrovsky,
Searchable symmetric encryption: improved definitions and efficient
constructions,
in \emph{Proc.\ CCS 2006}, pp.\ 79--88.

\bibitem{boneh2011fe}
D.~Boneh, A.~Sahai, B.~Waters,
Functional encryption: definitions and challenges,
in \emph{Proc.\ TCC 2011}, pp.\ 253--273.

\bibitem{pandey2012ppe}
O.~Pandey, Y.~Rouselakis,
Property preserving symmetric encryption,
in \emph{Proc.\ EUROCRYPT 2012}, pp.\ 375--391.

\bibitem{ateniese2005pre}
G.~Ateniese, K.~Fu, M.~Green, S.~Hohenberger,
Improved proxy re-encryption schemes with applications to secure
distributed storage,
\emph{ACM Trans.\ Inf.\ Syst.\ Secur.}\ \textbf{9}(1), 1--30 (2006).

\bibitem{libert2008pre}
B.~Libert, D.~Vergnaud,
Unidirectional chosen-ciphertext secure proxy re-encryption,
in \emph{Proc.\ Public Key Cryptography 2008}, pp.\ 360--379.

\bibitem{nunez2018umbral}
D.~Nu\~nez,
Umbral: a threshold proxy re-encryption scheme,
\emph{NuCypher whitepaper}, 2018.

\bibitem{shamir1979}
A.~Shamir,
How to share a secret,
\emph{Commun.\ ACM} \textbf{22}(11), 612--613 (1979).

\bibitem{rfc6962}
B.~Laurie, A.~Langley, E.~Kasper,
Certificate Transparency,
RFC 6962, IETF, 2013.

\bibitem{rfc8452}
S.~Gueron, A.~Langley, Y.~Lindell,
AES-GCM-SIV: nonce misuse-resistant authenticated encryption,
RFC 8452, IETF, 2019.

\bibitem{rfc5869}
H.~Krawczyk, P.~Eronen,
HMAC-based Extract-and-Expand Key Derivation Function (HKDF),
RFC 5869, IETF, 2010.

\bibitem{krawczyk2010hkdf}
H.~Krawczyk,
Cryptographic extraction and key derivation: the HKDF scheme,
in \emph{Proc.\ CRYPTO 2010}, pp.\ 631--648.

\bibitem{nist-sp800-38b}
M.~Dworkin,
Recommendation for block cipher modes of operation: the CMAC mode for
authentication,
NIST Special Publication 800-38B, 2005.

\bibitem{marlinspike2016signal}
M.~Marlinspike, T.~Perrin,
The Double Ratchet Algorithm,
Open Whisper Systems, 2016.

\bibitem{kohlweiss2023blueprints}
M.~Kohlweiss, A.~Lysyanskaya, A.~D.~Nguyen,
Privacy-preserving blueprints,
in \emph{Proc.\ EUROCRYPT 2023}, pp.\ 594--625.

\bibitem{camenisch2004clsigs}
J.~Camenisch, A.~Lysyanskaya,
Signature schemes and anonymous credentials from bilinear maps,
in \emph{Proc.\ CRYPTO 2004}, pp.\ 56--72.

\bibitem{camenisch2019mercurial}
E.~C.~Crites, A.~Lysyanskaya,
Mercurial signatures for delegatable anonymous credentials,
in \emph{Proc.\ CT-RSA 2019}, pp.\ 535--555.

\bibitem{garman2024mercurial}
C.~Garman, A.~Lysyanskaya, et al.,
Stronger privacy guarantees for delegatable anonymous credentials,
in \emph{Proc.\ ASIACRYPT 2024}.

\bibitem{lysyanskaya2002vrf}
A.~Lysyanskaya,
Unique signatures and verifiable random functions from the DH-DDH separation,
in \emph{Proc.\ CRYPTO 2002}, pp.\ 597--612.

\bibitem{camenisch2002accumulators}
J.~Camenisch, A.~Lysyanskaya,
Dynamic accumulators and application to efficient revocation of anonymous credentials,
in \emph{Proc.\ CRYPTO 2002}, pp.\ 61--76.

\bibitem{kemmoe2024accumulator}
V.~Y.~Kemmoe, A.~Lysyanskaya,
RSA-based dynamic accumulator without hashing into primes,
in \emph{Proc.\ ACM CCS 2024}.

\bibitem{chase2006sigs-of-knowledge}
M.~Chase, A.~Lysyanskaya,
Signatures of knowledge,
in \emph{Proc.\ CRYPTO 2006}, pp.\ 78--96.

\bibitem{bellare2024prfs}
M.~Bellare, A.~Lysyanskaya,
Symmetric and dual PRFs from standard assumptions: a generic validation of an HMAC assumption,
\emph{J.\ Cryptol.}\ \textbf{37}, 2024.

\bibitem{camenisch2005ecash}
J.~Camenisch, S.~Hohenberger, A.~Lysyanskaya,
Compact e-cash,
in \emph{Proc.\ EUROCRYPT 2005}, pp.\ 302--321.

\bibitem{goodrich2010authdata-generic}
M.~T.~Goodrich, C.~Papamanthou, R.~Tamassia, et al.,
Authenticated data structures, generically,
\emph{Theory of Computing Systems}, 2010.

\end{thebibliography}

\end{document}
